%- Package imports --------------------------------------------------------------------------------------------------------------------------------------
\usepackage[british, german]{babel}

\usepackage{amsmath} 
\usepackage{amsthm}
\usepackage{amstext}
\usepackage{amssymb}
\usepackage{amsfonts}
\usepackage{breqn}
\usepackage{csquotes}
\usepackage{graphicx} % Necessary to insert images
\usepackage{float} % Image positioning
\usepackage{caption}
\usepackage{enumitem} 
\usepackage{bm} % For the \bm command
\usepackage{setspace} % For line spacing etc.

\usepackage[a4paper, margin=32mm]{geometry}

% \usepackage{setspace}
% \setstretch{1.1}
% \usepackage{parskip}

\usepackage{fancyhdr}  
\fancyhf{} % Clear preset header and footer settings
\renewcommand\headrulewidth{0pt} % Do not display a line between header and content
\fancyfoot[CO, CE]{\thepage} % Display page number at the bottom center of each page
\pagestyle{fancy}

\usepackage[most]{tcolorbox}

\newtcolorbox{abstractframe}[1][]{
    enhanced,
    arc=0pt,
    outer arc=0pt,
    colback=white,
    boxrule=0.8pt,
    #1
}

\usepackage{titlesec} % Customize chapter headings
\titleformat{\chapter}
  {\normalfont\huge\bfseries}
  {\thechapter}{1em}{}

\usepackage{standalone}
  
\usepackage{tikz} % Graph drawing
\usetikzlibrary{graphs, angles, quotes, babel, matrix, decorations.pathreplacing, calc, positioning, backgrounds} 

\newlength\mtxrowsep   \setlength\mtxrowsep{1.5ex}
\newlength\mtxcolsep   \setlength\mtxcolsep{2\arraycolsep}

\tikzset{
  ams/.style={
    baseline=-.7ex,
    % every delimiter/.style={yshift=-1pt},
    every left delimiter/.style={xshift=4pt},
    every right delimiter/.style={xshift=-4pt},
    % every node/.style={inner sep=0pt},
  },
  ams matrix/.style={
    % inner sep=0pt,
    % column sep=\mtxcolsep,
    % row sep=\mtxrowsep,
    matrix of math nodes,
    baseline, 
    left delimiter={[},
    right delimiter={]}, 
    column sep=4pt
  },
  bmatrix/.style={
    ams,
    every matrix/.style={
      ams matrix,
      left delimiter={[},
      right delimiter={]},
    }
  },
}

\newcommand{\mynode} {
    \node[circle, draw=black, fill=black, inner sep=2pt, minimum size=3pt]
}

\tikzstyle{underbrace style}=[decorate,decoration={brace,raise=2mm,amplitude=3pt,mirror},color=gray]
\tikzstyle{underbrace text style}=[font=\tiny, below, pos=.5, yshift=-3mm]


\newcommand*{\boxcolor}{orange}
\makeatletter
\renewcommand{\boxed}[1]{
    \textcolor{\boxcolor}{%
        \tikz[baseline={([yshift=-1ex]current bounding box.center)}] 
        \node [rectangle, minimum width=1ex,rounded corners,draw, fill=red!20] {\normalcolor\m@th$\displaystyle#1$};
    }
}
\makeatother

 
%- Bibliography management ------------------------------------------------------------------------------------------------------------------------------
\usepackage[
    style=alphabetic, 
    backend=biber,
    maxbibnames=9
]{biblatex}

\usepackage{doi} % Displays DOI references as links

\addbibresource{ressources.bib}

\renewbibmacro{in:}{%
  \ifentrytype{article}{}{\printtext{\bibstring{in}\intitlepunct}}}
  
\DeclareFieldFormat[article,periodical]{volume}{\mkbibbold{#1}} % Displays the volume of a journal bold
\DeclareFieldFormat[inproceedings]{number}{\mkbibbold{#1}} % Displays the volume of proceedings bold
\DeclareFieldFormat{url}{Available at\addcolon\space\url{#1}} % Use "Available at: " as prefix for URLs

\DefineBibliographyExtras{british}{\def\finalandcomma{\addcomma}} % Enable Oxford comma

\DeclareSourcemap{ % Always use full journal name
  \maps[datatype=bibtex]{
    \map[overwrite=true]{
      \step[fieldsource=fjournal]
      \step[fieldset=journal, origfieldval]
    }
  }
}


%- Styles -----------------------------------------------------------------------------------------------------------------------------------------------
\theoremstyle{definition}
\newtheorem{definition}{Definition}[section]
\newtheorem{example}[definition]{Example}
\newtheorem{remark}[definition]{Remark}
\newtheorem{question}[definition]{Question}

\theoremstyle{plain}
\newtheorem{theorem}[definition]{Theorem}
\newtheorem{lemma}[definition]{Lemma}
\newtheorem{corollary}[definition]{Corollary}
\newtheorem{proposition}[definition]{Proposition}
\newtheorem{observation}[definition]{Observation}
\newtheorem{claim}{Claim}[definition]
\newtheorem{conjecture}[definition]{Conjecture}

\renewcommand{\labelenumi}{\rm{(\roman{enumi})}}

%- Custom math operators --------------------------------------------------------------------------------------------------------------------------------
\DeclareMathOperator{\tr}{tr}
\DeclareMathOperator{\dist}{dist}
\DeclareMathOperator{\vspan}{span}
\DeclareMathOperator{\rank}{rank}
\DeclareMathOperator*{\argmax}{arg\,max}
\DeclareMathOperator{\supp}{supp}

\newcommand{\N}{\mathbb{N}}
\newcommand{\Z}{\mathbb{Z}}
\newcommand{\Q}{\mathbb{Q}}
\newcommand{\R}{\mathbb{R}}
\newcommand{\C}{\mathbb{C}}

%- Images ------------------------------------------------------------------------------------------------------------------------------------------------
\graphicspath{ {./images/} }
